\frontmatterchapter{SAŽETAK}

U radu se razmatra problem autonomne interakcije mobilnog manipulatora s okolinom, na primjeru
otvaranja dvije vrste vrata, zakretnih i kliznih. Za mobilni manipulator KUKA KMR iiwa razvijen
je u simulacijskom okruženju cjelovit sustav koji objedinjuje percepciju, procjenu sile i
momenta te upravljanje u kontaktu. Modul percepcije iz slike određuje položaj i orijentaciju
kvake pomoću vizualnih oznaka, a sila i moment na vrhu alata procjenjuju se iz momenata u
zglobovima, bez zasebnog senzora sile. Tip vrata robot ne dobiva kao zadan podatak, nego ga
prepoznaje iz geometrije kvake.

Zadatak upravljanja riješen je dvama pristupima. Klasični pristup temeljen na modelu izvodi
cijeli lanac, od prilaska i hvata preko otvaranja do prolaska kroz vrata, oslanjajući se na
geometriju ograničenja izračunatu iz percepcije. Pristup temeljen na podržanom učenju zadatak
formulira kao Markovljev proces odlučivanja s varijabilnom impedancijom u prostoru akcije, pri
čemu politika u svakom koraku bira krutost regulatora po svakoj osi, a geometriju ograničenja
mora otkriti iz vlastitih opažanja sile i pomaka.

Oba su pristupa izmjerena na istom robotu, istoj sceni i istim mjerilima uspješnosti. Kod
kliznih vrata oba dosežu puni cilj u svakom pokušaju. Kod zakretnih vrata klasični pristup
pouzdano otvara vrata do 50\textdegree{}, ali preko otprilike 57\textdegree{} ne može jer kvaka
klizne iz prihvatnice, dok naučena politika vrata otvara preko 90\textdegree{} u 96 \% epizoda,
uz bitno niže kontaktne sile. Zaključak je da je eksplicitno poznavanje geometrije dovoljno kad
je geometrija jednostavna i dobro procijenjena, a gubi prednost čim zadatak zahtijeva prilagodbu
krutosti tijekom same interakcije.

\vspace{8mm}
\noindent\textbf{KLJUČNE RIJEČI:} mobilni manipulator, fizička interakcija, procjena
svojstava okoline, impedancijsko upravljanje, podržano učenje, simulacija.
\clearpage